\chapter{组合数学中的线性代数方法}
\label{ch:250526}

\ifdefined\VideoNotesTestMode
本章整理课程视频 Lec16 的课堂内容。三个录制分段按照给定顺序合并为一条连续时间轴；以下时间均指合并后的视频时间。参考 PDF 只用于核对术语，不作为本章的知识来源。
\fi

\section{欧氏空间中的距离集}

这一部分的共同思想是把距离条件写成内积、矩阵或多项式条件，再用秩或维数限制点集大小。

\subsection{两两等距的点集}

\begin{theorem}
在 $\mathbb{R}^d$ 中，两两距离均为 $1$ 的点至多有 $d+1$ 个。

\ifdefined\VideoNotesTestMode
\textbf{视频来源：Lec16，视频时间：00:08:50--00:27:19。}
\fi
\end{theorem}

\begin{proof}
反设有 $d+2$ 个这样的点。平移不改变点间距离，故可令其中一点为 $v_0=0$，其余点为 $v_1,\dots,v_{d+1}$，并置
\[
M=\begin{bmatrix}v_1&v_2&\cdots&v_{d+1}\end{bmatrix}
  \in\mathbb{R}^{d\times(d+1)}.
\]
由 $\lVert v_i\rVert=1$ 以及 $\lVert v_i-v_j\rVert=1$，对 $i\ne j$ 有
\[
2\langle v_i,v_j\rangle
=\lVert v_i\rVert^2+\lVert v_j\rVert^2-\lVert v_i-v_j\rVert^2=1.
\]
因此 Gram 矩阵满足
\[
M^TM=\frac12\bigl(I_{d+1}+\mathbf{1}\mathbf{1}^T\bigr).
\]
矩阵 $\mathbf{1}\mathbf{1}^T$ 在 $\mathbf{1}$ 方向上的特征值为 $d+1$，在其正交补上的特征值为 $0$，所以 $M^TM$ 的全部特征值均非零。于是
\[
\operatorname{rank}(M^TM)=d+1,
\]
但另一方面 $\operatorname{rank}(M^TM)\leq\operatorname{rank}(M)\leq d$，矛盾。

\ifdefined\VideoNotesTestMode
\textbf{视频来源：Lec16，视频时间：00:11:06--00:27:19。}
\fi
\end{proof}

这个证明体现了 Gram 矩阵的作用：点的长度与夹角都记录在 $M^TM$ 中，几何上的距离限制因此转化为矩阵的秩限制。正单纯形给出 $d+1$ 个两两等距点，所以该上界是紧的。

\ifdefined\VideoNotesTestMode
\textbf{视频来源：Lec16，视频时间：00:08:50--00:10:22。}
\fi

\subsection{两两为奇整数距离的点集}

\begin{theorem}
若 $X\subset\mathbb{R}^d$ 中任意两点之间的距离都是奇整数，则
\[
|X|\leq d+2.
\]

\ifdefined\VideoNotesTestMode
\textbf{视频来源：Lec16，视频时间：00:29:00--00:44:52。}
\fi
\end{theorem}

\begin{proof}
设 $|X|=n+1$。平移后写成 $X=\{0,v_1,\dots,v_n\}$，并令
\[
M=\begin{bmatrix}v_1&\cdots&v_n\end{bmatrix},
\qquad B=2M^TM.
\]
因为 $\lVert v_i\rVert$ 和 $\lVert v_i-v_j\rVert$ 都是奇整数，$B$ 的对角元
$2\lVert v_i\rVert^2$ 为偶数，而当 $i\ne j$ 时
\[
B_{ij}=\lVert v_i\rVert^2+\lVert v_j\rVert^2-
\lVert v_i-v_j\rVert^2
\]
为奇数。因此在 $\mathbb{F}_2$ 上
\[
B\equiv J_n-I_n,
\]
其中 $J_n$ 是全 $1$ 矩阵。

若 $n$ 为偶数，则 $J_n-I_n$ 在 $\mathbb{F}_2$ 上可逆。课堂从行列式展开解释了这一点：模 $2$ 后符号消失，非零项恰好对应没有不动点的排列，故行列式的奇偶性等于错排数的奇偶性；偶数阶错排数为奇数。于是 $\operatorname{rank}(B)=n$，从而 $n\leq d$。

若 $n$ 为奇数，则删去一个 $v_i$ 后得到偶数阶的同型主子矩阵，故 $\operatorname{rank}(B)\geq n-1$，从而 $n-1\leq d$。两种情形合并即得 $|X|=n+1\leq d+2$。

\ifdefined\VideoNotesTestMode
\textbf{视频来源：Lec16，视频时间：00:29:00--00:44:52。}
\fi
\end{proof}

这里先把整数矩阵模 $2$，是因为若其模 $2$ 的行列式非零，则原整数行列式必为奇数，因而在 $\mathbb{R}$ 上也非零。这个步骤把度量问题连接到了有限域上的线性代数。

\textbf{待核实：}课堂随后提到还可继续研究模 $4$、模 $8$、模 $16$ 的信息，并口头提及一个与模 $16$ 余数 $14$ 有关的情形；音频与板书不足以唯一确定其精确陈述，因此这里不写成结论。

\ifdefined\VideoNotesTestMode
\textbf{视频来源：Lec16，视频时间：00:45:00--00:46:06。}
\fi

\subsection{两距离点集}

令 $f(d)$ 表示 $\mathbb{R}^d$ 中两两距离恰好至多取两个值的有限点集的最大大小。

\ifdefined\VideoNotesTestMode
\textbf{视频来源：Lec16，视频时间：00:46:14--00:49:08。}
\fi

\begin{example}
取 $\mathbb{R}^{d+1}$ 中所有恰有两个坐标为 $1$、其余坐标为 $0$ 的点。任意两个不同点的距离为 $\sqrt{2}$ 或 $2$。
按照常规思路，我们取的是 $\mathbb{R}^{d+1}$，因此我们得到的结果是关于 $f(d+1)$ 的。不过注意到，这些点都位于 $d$ 维仿射超平面
\[
x_1+\cdots+x_{d+1}=2
\]
中，故实际上这是 $\mathbb R^d$ 的情况，那么
\[
f(d)\geq\binom{d+1}{2}=\frac{d(d+1)}2.
\]

\ifdefined\VideoNotesTestMode
\textbf{视频来源：Lec16，视频时间：00:47:29--00:49:08。}
\fi
\end{example}


\section{线性无关}

两距离集上界所用的基本技巧，是构造一族在给定点集上具有简单取值模式的函数。

\begin{lemma}
设 $g_1,\dots,g_m$ 是函数，$a_1,\dots,a_m$ 是其定义域中的点。若评值矩阵
\[
G=\bigl(g_j(a_i)\bigr)_{1\leq i,j\leq m}
\]
可逆，则 $g_1,\dots,g_m$ 线性无关。特别地，对角线上非零而其余位置为零的评值矩阵，以及对角线上非零的三角评值矩阵，都满足这一条件。

\ifdefined\VideoNotesTestMode
\textbf{视频来源：Lec16，视频时间：01:01:30--01:17:59。}
\fi
\end{lemma}

\begin{proof}
若 $\sum_{j=1}^m\lambda_jg_j=0$，依次代入 $a_1,\dots,a_m$，便得到
$G\lambda=0$。由 $G$ 可逆可知 $\lambda=0$。
\end{proof}

这条引理连接了两个维数观点：一方面，评值矩阵证明所构造的函数至少张成一个 $m$ 维空间；另一方面，若这些函数都落在一个显式的 $D$ 维多项式空间中，就必有 $m\leq D$。

\begin{theorem}
两距离集满足
\[
f(d)\leq\frac{(d+1)(d+4)}2.
\]

\ifdefined\VideoNotesTestMode
\textbf{视频来源：Lec16，视频时间：00:49:04--01:11:22。}
\fi
\end{theorem}

\begin{proof}
设 $a_1,\dots,a_m\in\mathbb{R}^d$ 的两种非零距离为 $\delta_1,\delta_2$。对每个 $k$ 定义
\[
g_k(x)=\bigl(\lVert x-a_k\rVert^2-\delta_1^2\bigr)
       \bigl(\lVert x-a_k\rVert^2-\delta_2^2\bigr).
\]
若 $i\ne k$，则 $\lVert a_i-a_k\rVert$ 等于 $\delta_1$ 或 $\delta_2$，所以 $g_k(a_i)=0$；而
\[
g_k(a_k)=\delta_1^2\delta_2^2\ne0.
\]
评值矩阵是非零对角矩阵，故 $g_1,\dots,g_m$ 线性无关。

展开 $g_k$ 可知它属于
\[
\operatorname{span}\left\{
1,\ x_i,\ x_ix_j\ (i\leq j),\
x_i\sum_{j=1}^d x_j^2,\
\left(\sum_{i=1}^d x_i^2\right)^2
\right\}.
\]
这里依次有 $1,d,\binom{d+1}{2},d,1$ 个生成元，所以该空间的维数至多
\[
1+d+\binom{d+1}{2}+d+1
=\frac{(d+1)(d+4)}2.
\]
线性无关函数的个数不能超过所在空间的维数，结论得证。

\ifdefined\VideoNotesTestMode
\textbf{视频来源：Lec16，视频时间：00:49:04--01:11:22。}
\fi
\end{proof}

\ifdefined\VideoNotesTestMode
第一段视频在写出 $g_k$ 后进入课间，第二段视频从评值矩阵继续该证明。合并后的连续时间轴在此处没有把一个数学结论拆成两个内容单元。
\fi

\begin{theorem}
两距离集还有更强的上界
\[
f(d)\leq\frac{(d+1)(d+2)}2.
\]

\ifdefined\VideoNotesTestMode
\textbf{视频来源：Lec16，视频时间：01:17:59--01:40:38。}
\fi
\end{theorem}

\begin{proof}
沿用上面的 $a_k$ 与 $g_k$。证明的关键是：不仅 $g_1,\dots,g_m$ 线性无关，而且
\[
g_1,\dots,g_m,1,x_1,\dots,x_d
\]
也线性无关。

设存在恒等式
\[
\sum_{k=1}^m\lambda_kg_k(x)+\sum_{p=1}^d\mu_px_p+\tau=0.
\]
比较四次齐次项可得 $\sum_k\lambda_k=0$。再对不同的 $p,q$ 求混合二阶偏导。由于
\[
\frac{\partial^2 g_k}{\partial x_p\partial x_q}
=8(x_p-a_{kp})(x_q-a_{kq})\qquad(p\ne q),
\]
比较所得二次多项式的系数
\[
   \sum_{i\in[m]}\lambda_i(x_p-a_{ip})(x_q-a_{iq})=x_px_q\sum_{i=1}^m\lambda_i-x_p\sum_{i=1}^m\lambda_ia_{iq}-x_q\sum_{i=1}^m\lambda_ia_{ip}+\sum_{i=1}^m\lambda_ia_{ip}a_{iq}=0
\]
可推出
\[
\sum_{k=1}^m\lambda_k=\sum_{k=1}^m\lambda_ka_{kp}=0
\qquad(1\leq p\leq d).
\]

把原恒等式代入各点 $a_i$。记 $c=\delta_1^2\delta_2^2>0$，则
\[
c\lambda_i+\sum_{p=1}^d\mu_pa_{ip}+\tau=0
\qquad(1\leq i\leq m).
\]
两边乘以 $\lambda_i$ 后对 $i$ 求和，并使用上面得到的
$\sum_{i=1}^m\lambda_i=0$ 与 $\sum_{i=1}^m\lambda_i a_{ip}=0$，便有
\[
c\sum_{i=1}^m\lambda_i^2=0.
\]
因为 $c>0$，所以所有 $\lambda_i=0$；随后一次多项式
$\sum_p\mu_px_p+\tau$ 恒为零，故所有 $\mu_p$ 和 $\tau$ 也为零。

这 $m+d+1$ 个多项式仍落在上一证明中的空间里，因此
\[
m+d+1\leq\frac{(d+1)(d+4)}2,
\]
整理即得 $m\leq (d+1)(d+2)/2$。

\ifdefined\VideoNotesTestMode
\textbf{视频来源：Lec16，视频时间：01:17:59--01:40:38。}
\fi
\end{proof}

这个改进没有缩小多项式空间，而是向已知线性无关组中再加入 $d+1$ 个低次多项式。求导的作用是消去这些低次项并读取高次系数；最后的平方和则把若干线性条件合并成 $\lambda_i=0$ 的强结论。

\section{石头剪刀布的多项式编码}

\begin{example}
把石头、剪刀、布编码为 $\mathbb{F}_3$ 中的三个元素，并约定相差 $1$ 的方向表示“胜过”，例如 0 胜过 2。设每位参与者预先提交一个长度为 $n$ 的序列
$a=(a_1,\dots,a_n)\in\mathbb{F}_3^n$。
令 $R(n)$ 为满足以下性质的序列族的最大规模：任意两个不同参与者 A, B, 至少有一局 A 胜过 B，且至少有一局 B 胜过 A。则
\[
\binom{n}{\lfloor n/2\rfloor}\leq R(n)\leq2^n.
\]

\ifdefined\VideoNotesTestMode
\textbf{视频来源：Lec16，视频时间：01:41:28--01:59:22。}
\fi
\end{example}

下界可取所有重量为 $\lfloor n/2\rfloor$ 的 $0$--$1$ 序列。两个不同的等重序列必同时出现 $0$ 变 $1$ 和 $1$ 变 $0$ 的坐标，因而两个方向的取胜条件都能实现；其规模约为 $2^n/\sqrt{n}$。

上界使用同一评值思想。对每个序列 $a$ 构造
\[
h_a(x)=\prod_{i=1}^n(x_i-a_i-1)\in\mathbb F_3[X_1,\dots,X_n].
\]
有 $h_a(a)=(-1)^n\ne0$；对任何另一序列 $b$，有向取胜条件保证某个因子为零，因此 $h_a(b)=0$。故这些 $\{h_a\}_{a\in\mathbb{F}_3^n}$ 线性无关。每个变量在 $h_a$ 中的次数至多为 $1$，而所有多重线性单项式
\[
\prod_{i\in S}x_i\qquad(S\subseteq[n])
\]
只有 $2^n$ 个，所以 $R(n)\leq2^n$。

\section{Joints 问题}

\begin{definition}[Joint]
给定 $\mathbb{R}^d$ 中的一族直线，若一点同时位于其中 $d$ 条直线上，且这 $d$ 条直线的方向向量线性无关，则称该点为一个 joint。

\ifdefined\VideoNotesTestMode
\textbf{视频来源：Lec16，视频时间：02:00:00--02:03:05。}
\fi
\end{definition}

我们记 $d$ 维空间中 $N$ 条直线能相交确定的 joints 数量的上限为 $J_d(N)$. 
显然有 $J_2(N)={N\choose 2}$，因为 $J_d$ 公共上界就是 ${N\choose 2}$; 而且在二维空间，可以构造出这么多个 2-joints。

给出 $\mathbb{R}^3$ 中的一种构造：取三个坐标方向、每个方向各约 $k^2$ 条网格线，共有 $N\asymp3k^2$ 条线，产生 $k^3\asymp N^{3/2}$ 个 joints。
这说明了一个下界 \(\Omega(N^{1.5}).\) 下面证明这同时也是上界。

\begin{theorem}
$\mathbb{R}^3$ 中 $N$ 条直线所确定的 joints 数量为 \(O(N^{1.5}).\)

\ifdefined\VideoNotesTestMode
\textbf{视频来源：Lec16，视频时间：02:00:00--02:30:58。}
\fi
\end{theorem}

\begin{proof}
反设对于任何常数 $C$, 总有足够大的 $N$，使得存在一个大小为 $N$ 的直线族 $L$, 
其 joints 数量 $M>CN^{1.5}$。下面任取一个 $C$，记对应的直线族和 joints 集合为 $L,J$，$|L|=N,|J|=M$.

反复删除当前所含 joints 个数少于 $M/(2N)$ 个的直线，
并同时删除这些直线上的当前 joints 
(注意虽然这里删除了线和 joints, 我们仍然用 $N,M$ 记录删除前的线数和 joint 数).
整个过程删除的 joints 少于 $M/2$，
所以留下\textbf{非空的}直线族 $L^*$ 与 joint 集 $J^*$；
其中每条直线都含有至少 $M/(2N)$ 个删除后留下的 joints。

我们通过多项式来导出矛盾，这里需要一个在 $J^*$ 上取值为 $0$ 的多项式。

三变量总次数不超过 $r$ 的多项式空间维数为
\[
\binom{r+3}{3}.
\]
这是因为单项式 $x_1^{e_1}x_2^{e_2}x_3^{e_3}$ 对应满足
$e_1+e_2+e_3\leq r$ 的非负整数解（插板法）。
取 $r=\sqrt[3]{6M}$，使上述维数 $\binom{r+3}{3}>\frac{r^3}{6}=M$，
那么存在一个非零多项式 $g$，次数为 $r$，
且对于所有点 $A\in J^*$，$g(A)=0$ (待定系数法)。

把 $g$ 限制到任意一条直线 $\ell:\mathbf x=\mathbf a+t\mathbf v\in L^*$，得到关于 $t$ 的一元多项式，
次数不超过 $r$。由于 $N<(M/C)^{2/3}$，该直线上的零点数 
$|\ell\cap P^*|\geq \frac{M}{2N}>\sqrt[3]{MC^2/8}$. 
如果我们取 $C=8$, 那么 $|\ell\cap P^*|>r$, 所以多项式 $g$ 在整条直线上取值恒为 $0$。
由于刚才任选的直线，因此对于所有直线族中的直线来说，$g$ 在这些直线上取值都是 $0$.

考虑 joint $P\in J^*$，$g$ 沿三条线性无关方向都恒为零，故三个方向导数都为零。方向导数等于方向向量与梯度的内积；三个方向张成 $\mathbb{R}^3$，因此 $\nabla P=0$。于是每个一阶偏导数也在所有留下的 joints 上为零。

每个非零一阶偏导数的次数低于 $r$, 因此 $\frac{\partial g}{\partial x_i}$ 也是一个多项式，满足在 $J^*$ 上取值为 $0$, 且次数至多 $r$。
重复“直线上零点多于次数”的论证，可知它在每条留下的直线上恒为零，
继而它的全部一阶偏导数也在 joints 上为零。
如此迭代，$g$ 的所有阶、任何一个偏导数都在 $J^*$ 上为零。
假设 $g$ 存在非零最高次项，那么总有一个高阶偏导数在全空间是常数。因此 $g=0$，矛盾。

\ifdefined\VideoNotesTestMode
\textbf{视频来源：Lec16，视频时间：02:06:40--02:30:58。}
\fi
\end{proof}

证明中的三项基础事实分别是：维数大于约束数时齐次线性方程组有非零解；一元多项式的零点数不超过其次数；三个独立方向上的方向导数为零等价于梯度为零。把它们串联起来，就能让“在有限点集上消失”逐步升级为“所有阶导数都消失”。

对于 $d$ 维的情况，joint 个数为
\[
O\bigl(N^{d/(d-1)}\bigr),
\]
同样下界由网格构造给出。

\ifdefined\VideoNotesTestMode
\textbf{视频来源：Lec16，视频时间：02:03:05--02:05:00。}
\fi

\section{Finite Kakeya Problem}

\begin{definition}[有限 Kakeya 集]
设 \(K\subseteq\mathbb{F}_q^d\)。若对每个非零方向
\(\mathbf v\in\mathbb{F}_q^d\setminus\{0\}\)，都存在
\(\mathbf a\in\mathbb{F}_q^d\)，使得直线
\[
\{\mathbf a+t\mathbf v:t\in\mathbb{F}_q\}\subseteq K,
\]
则称 \(K\) 为有限 Kakeya 集。

\ifdefined\VideoNotesTestMode
\textbf{视频来源：Lec16，视频时间：02:31:50--02:35:30。}
\fi
\end{definition}

\begin{theorem}[有限 Kakeya 定理(Proved in 2009)]
对每个固定维数 \(d\)，存在只依赖于 \(d\) 的常数
\(c_d>0\)，使得对任意有限域 \(\mathbb{F}_q\) 以及任意有限
Kakeya 集 \(K\subseteq\mathbb{F}_q^d\)，都有
\[
|K|\geq c_dq^d.
\]
\end{theorem}

它与 Joints 证明的联系仍是低次数多项式：若 $K$ 太小，维数计数会给出一个在 $K$ 上消失的非零低次多项式；把它限制到 $K$ 所包含的各方向直线上，再利用一元多项式的零点性质，就能把信息传递到最高次齐次部分。

\begin{proof}
反设存在 $d$, 对任意 $c_d$ 存在有限域 \(\mathbb{F}_q\) 以及有限
Kakeya 集 \(K\subseteq\mathbb{F}_q^d\), 使得 $|K|< c_dq^d$. 
下面任取 $c_d$，然后取对应的有限域 \(\mathbb{F}_q\) 和有限 Kakeya 集 \(K\subseteq\mathbb{F}_q^d\)
满足 $|K|=m< c_dq^d$. 

我们需要 $K$ 上面取值为零的多项式 $g\in\mathbb F_q[X_1,\dots,X_d]$. 
对于 $r=q\left(d!c_d\right)^{1/d}$, 由于 ${r+d\choose d}>\frac{r^d}{d!}=c_dq^d>m$，一定存在一个
次数最高为 $r$ 的非零多项式 $g$ 在 $K$ 上取值为 0.

由于 $c_d$ 任取，不妨让 $c_d<\frac{1}{d!}$, 此时由于直线 
$\{\mathbf a+t\mathbf v:t\in \mathbb F_q\}$ 
上有 $q$ 个点，$g$ 在直线上的限制是一个低于 $q$ 次的一元多项式，故为零多项式。

可以计算得到 $g(\mathbf a+t\mathbf v)$ 展开式中，关于 $t$ 的最高次项为 $g_r(\mathbf v)t^r$, 
其中 $g_r$ 为 $g$ 中次数为 $r$ 的齐次多项式。Kakeya 集合的性质可表述为：对于任意 $\mathbf v$,
$g(\mathbf a+t\mathbf v)$ 关于 $t$ 是零多项式，
得到 $g_r(\mathbf v)t^r=0$ 对于任意 $\mathbf v\in\mathbb F_q^d\setminus\{0\}$ 都成立。
由齐次性质 $g_r(0)=0$. 但是一个非零多项式的最高次不可能是零多项式，矛盾。

\end{proof}

\ifdefined\VideoNotesTestMode
\textbf{视频来源：Lec16，视频时间：02:31:50--02:38:02。}
\fi
